% JAIR's reproducibility checklist, verbatim from their template with the
% answers filled in. MANDATORY: "Submissions without a completed
% reproducibility checklist will be desk rejected." Used only in review; it
% comes out of the camera-ready.
%
% Two answers are deliberately "partially" and two blocks are "NA". Do not
% round them up to "yes" to make the page look better. A checklist that
% overstates is worse than one that admits a gap, and it would be a strange
% paper to lie on this particular form.

\section{Reproducibility Checklist for JAIR}

Select the answers that apply to your research -- one per item.

\subsection*{All articles:}

\begin{enumerate}
    \item All claims investigated in this work are clearly stated.
    \textbf{[yes]}
    \item Clear explanations are given how the work reported substantiates the
    claims. \textbf{[yes]}
    \item Limitations or technical assumptions are stated clearly and
    explicitly. \textbf{[yes]}
    \item Conceptual outlines and/or pseudo-code descriptions of the AI methods
    introduced in this work are provided, and important implementation details
    are discussed. \textbf{[yes]}
    \item
    Motivation is provided for all design choices, including algorithms,
    implementation choices, parameters, data sets and experimental protocols
    beyond metrics. \textbf{[yes]}
\end{enumerate}

\subsection*{Articles containing theoretical contributions:}
Does this paper make theoretical contributions? \textbf{[yes]}

If yes, please complete the list below.

\begin{enumerate}
    \item All assumptions and restrictions are stated clearly and formally.
    \textbf{[yes]}
    \item All novel claims are stated formally (e.g., in theorem statements).
    \textbf{[yes]}
    \item Proofs of all non-trivial claims are provided in sufficient detail to
    permit verification by readers with a reasonable degree of expertise (e.g.,
    that expected from a PhD candidate in the same area of AI).
    \textbf{[yes]}
    \item
    Complex formalism, such as definitions or proofs, is motivated and
    explained clearly. \textbf{[yes]}
    \item
    The use of mathematical notation and formalism serves the purpose of
    enhancing clarity and precision; gratuitous use of mathematical formalism
    (i.e., use that does not enhance clarity or precision) is avoided.
    \textbf{[yes]}
    \item
    Appropriate citations are given for all non-trivial theoretical tools and
    techniques. \textbf{[yes]}
\end{enumerate}

\subsection*{Articles reporting on computational experiments:}
Does this paper include computational experiments? \textbf{[yes]}

If yes, please complete the list below.
\begin{enumerate}
    \item
    All source code required for conducting experiments is included in an
    online appendix or will be made publicly available upon publication of the
    paper. The online appendix follows best practices for source code
    readability and documentation as well as for long-term accessibility.
    \textbf{[yes]}
    \item The source code comes with a license that
    allows free usage for reproducibility purposes. \textbf{[yes]}
    \item The source code comes with a license that
    allows free usage for research purposes in general. \textbf{[yes]}
    \item
    Raw, unaggregated data from all experiments is included in an online
    appendix or will be made publicly available upon publication of the paper.
    The online appendix follows best practices for long-term accessibility.
    \textbf{[yes]}
    \item The unaggregated data comes with a license that
    allows free usage for reproducibility purposes. \textbf{[yes]}
    \item The unaggregated data comes with a license that
    allows free usage for research purposes in general. \textbf{[yes]}
    \item If an algorithm depends on randomness, then the method used for
    generating random numbers and for setting seeds is described in a way
    sufficient to allow replication of results. \textbf{[yes]}
    \item The execution environment for experiments, the computing
    infrastructure (hardware and software) used for running them, is described,
    including GPU/CPU makes and models; amount of memory (cache and RAM); make
    and version of operating system; names and versions of relevant software
    libraries and frameworks. \textbf{[partially]}
    \item
    The evaluation metrics used in experiments are clearly explained and their
    choice is explicitly motivated. \textbf{[yes]}
    \item
    The number of algorithm runs used to compute each result is reported.
    \textbf{[yes]}
    \item
    Reported results have not been ``cherry-picked'' by silently ignoring
    unsuccessful or unsatisfactory experiments. \textbf{[yes]}
    \item
    Analysis of results goes beyond single-dimensional summaries of performance
    (e.g., average, median) to include measures of variation, confidence, or
    other distributional information. \textbf{[yes]}
    \item
    All (hyper-) parameter settings for
    the algorithms/methods used in experiments have been reported, along with
    the rationale or method for determining them. \textbf{[yes]}
    \item
    The number and range of (hyper-) parameter settings explored prior to
    conducting final experiments have been indicated, along with the effort
    spent on (hyper-) parameter optimisation. \textbf{[partially]}
    \item
    Appropriately chosen statistical hypothesis tests are used to establish
    statistical significance in the presence of noise effects.
    \textbf{[yes]}
\end{enumerate}


\subsection*{Articles using data sets:}
Does this work rely on one or more data sets (possibly obtained from a
benchmark generator or similar software artifact)? \textbf{[yes]}

If yes, please complete the list below.
\begin{enumerate}
    \item
    All newly introduced data sets
    are included in an online appendix
    or will be made publicly available upon publication of the paper.
    The online appendix follows best practices for long-term accessibility with
    a license that allows free usage for research purposes. \textbf{[NA]}
    \item The newly introduced data set comes with a license that
    allows free usage for reproducibility purposes. \textbf{[NA]}
    \item The newly introduced data set comes with a license that
    allows free usage for research purposes in general. \textbf{[NA]}
    \item All data sets drawn from the literature or other public sources
    (potentially including authors' own previously published work) are
    accompanied by appropriate citations. \textbf{[yes]}
    \item All data sets drawn from the existing literature (potentially
    including authors' own previously published work) are publicly available.
    \textbf{[yes]}
    \item All new data sets and data sets that are not publicly available are
    described in detail, including relevant statistics, the data collection
    process and annotation process if relevant. \textbf{[NA]}
    \item
    All methods used for preprocessing, augmenting, batching or splitting data
    sets (e.g., in the context of hold-out or cross-validation)
    are described in detail. \textbf{[yes]}
\end{enumerate}

\subsection*{Explanations on any of the answers above (optional):}

\emph{Computational experiments, items 1 and 4: which clause we satisfy.} We
provide no online appendix. Both items are met by their second clause: the
source code and the raw per-cell result files are already public, at
\url{\repourl}, rather than deposited with this paper. That is deliberate.
What Appendix~\ref{app:prereg} asks a reader to check is the \emph{order} in
which things were committed, and only a repository with its commit graph
intact supports that: the reader runs \texttt{git merge-base} and gets an
answer. A static archive would turn the same claim back into something taken
on trust, which is the one thing this paper is trying not to ask for.

\emph{Data sets, items 1--3 and 6.} This work introduces no new data set. All
four training tasks and both downstream benchmarks are public and cited.

\emph{Computational experiments, item 8: partially.} The GPU make and model
and the per-cell memory footprint are reported in the Reproducibility
Statement, as is the aggregate compute. Operating system
version and pinned library versions are recorded in the released environment
files rather than in the paper.

\emph{Computational experiments, item 14: partially.} The regularisation sweep
reports its full grid and range. We do not report a systematic accounting of
settings explored before the final experiments, because most were not
hyperparameter search: they were superseded runs discarded after defects were
found in the scorers and in one data loader, and those are described in
\S\ref{sec:discussion} rather than summarised as tuning effort.

\emph{Item 11, on cherry-picking.} Every empirical claim in this paper is
governed by a pre-registration committed to version control before the compute
it governs was dispatched, with all outcome branches written in advance. The
registrations are included unedited and the commit ordering is verifiable from
the released history with \texttt{git merge-base}. Several of our own earlier
claims did not survive these tests and are reported as withdrawn rather than
removed.
